\documentclass{article}
\usepackage{amsmath,amssymb}
\usepackage[a4paper,margin=2.5cm]{geometry}
\usepackage{graphicx}
\usepackage{hyperref}

\title{Otimização de Portfólio de Energia com Previsão de PLD e Heurística em Grafos para Seleção de Trades}
\author{Proposta de Trabalho de Conclusão de Curso}
\date{2025}

\begin{document}
\maketitle

% ============================================================
\section{Introdução}
% ============================================================

Esta proposta trata de um problema muito comum em finanças, mas aplicado à energia elétrica:
\textbf{decidir compras e vendas (trades) para melhorar o resultado financeiro}, usando
\textbf{previsões de preço} e \textbf{controle de risco}.

No Brasil, o preço de referência do mercado de curto prazo de energia é o \textbf{PLD}
(Preço de Liquidação de Diferenças), em R\$/MWh. Para este trabalho, pense no PLD como:
\begin{quote}
\emph{``o preço do mercado naquele mês (ou semana), que ninguém sabe exatamente qual será no futuro''}.
\end{quote}

Vamos trabalhar com uma
visão simplificada e prática:
\begin{itemize}
    \item existe um portfólio com \textbf{contratos de compra e venda} (quantidade em MWh e preço em R\$/MWh);
    \item existe uma \textbf{produção} (por exemplo, de uma usina), que pode ser conhecida ou incerta;
    \item a \textbf{diferença} entre o que a empresa tem (produção + compras) e o que ela prometeu entregar (vendas)
    é liquidada ao PLD.
\end{itemize}

O que este trabalho acrescenta é uma \textbf{camada de decisão}:
\begin{quote}
\emph{``quais novos contratos devo comprar/vender agora para ter melhor resultado e menor risco no futuro?''}
\end{quote}

% ============================================================
\section{Dados}
% ============================================================

Vamos assumir que temos:

\subsection{PLD histórico (passado)}
Uma planilha com valores passados do PLD por submercado (por exemplo: Sudeste, Sul, Nordeste e Norte),
com uma data por mês e uma coluna de PLD para cada submercado.

\subsection{Projeções de PLD via NEWAVE (futuro)}
Podemos usar \textbf{séries sintéticas de CMO do NEWAVE} para compor projeções
de PLD quando entramos em um período futuro.
A ideia essencial é:
\begin{itemize}
    \item o NEWAVE gera muitos futuros possíveis;
    \item cada futuro é um \textbf{cenário} com uma trajetória de custo (CMO), que usaremos como proxy de preço futuro;
    \item São usados \textbf{2000 cenários} como saida.
\end{itemize}

\subsection{Carteira de contratos (compras e vendas)}
Cada contrato tem:
\begin{itemize}
    \item \textbf{quantidade} (MWh);
    \item \textbf{preço} (R\$/MWh);
    \item \textbf{período} (ex.: mês);
    \item \textbf{submercado}.
\end{itemize}

\subsection{Produção (opcionalmente incerta)}
A produção pode também ser incerteza (por exemplo, eólicas, usando distribuições por mês).
No trabalho, isso pode ser tratado de duas formas:
\begin{itemize}
    \item \textbf{versão simples}: produção futura conhecida (ou uma previsão única);
    \item \textbf{versão completa}: produção futura por cenário.
\end{itemize}

% ============================================================
\section{Conceito-chave: como a empresa ``ganha ou perde'' dinheiro}
% ============================================================

Esta é a parte mais importante para entender: \textbf{o resultado financeiro tem duas partes}.

\subsection{Parte 1: contratos a preço fixo}
Se eu \textbf{vendo} $Q$ MWh por um preço fixo $K$ (R\$/MWh), eu recebo:
\[
\text{Receita fixa} = K \cdot Q.
\]
Se eu \textbf{compro} $Q$ MWh por um preço fixo $K$, eu pago:
\[
\text{Custo fixo} = K \cdot Q.
\]

\subsection{Parte 2: ajuste ao PLD (o ``resto'' vai para o mercado de curto prazo)}
No mês, a empresa tem:
\begin{itemize}
    \item \textbf{energia disponível}: produção + compras;
    \item \textbf{energia comprometida}: vendas.
\end{itemize}
A diferença (que pode ser positiva ou negativa) é liquidada ao PLD:
\[
\text{Exposição ao PLD} = (\text{produção} + \text{compras} - \text{vendas}).
\]
\begin{itemize}
    \item Se a exposição é positiva, sobrou energia: em termos simples, ``vende no PLD''.
    \item Se a exposição é negativa, faltou energia: em termos simples, ``compra no PLD''.
\end{itemize}

\subsection{Exemplo de 1 mês (didático)}
Suponha:
\begin{itemize}
    \item produção = 90 MWh;
    \item venda fixa = 100 MWh a 200 R\$/MWh;
    \item PLD do mês = 300 R\$/MWh.
\end{itemize}
A empresa recebe no contrato: $200 \cdot 100 = 20.000$.
Mas falta energia: $90 - 100 = -10$ MWh, então ela compra 10 MWh ao PLD:
$-10 \cdot 300 = -3.000$.
Resultado do mês:
\[
20.000 - 3.000 = 17.000.
\]
Se o PLD fosse menor, a perda por faltar energia seria menor. Por isso o PLD é a grande incerteza.


% ============================================================
\section{Modelo de otimização}
% ============================================================

Agora vem o ``coração'': um modelo matemático que escolhe trades.

\subsection{O que é um ``trade'' nesse contexto?}
Um trade aqui é uma oportunidade de contratar energia (compra ou venda) para um submercado $s$
e um mês $t$, com um preço fixo conhecido no momento da decisão.

Exemplos:
\begin{itemize}
    \item ``comprar 50 MWh para o submercado SE em jul/2025 a 180 R\$/MWh'';
    \item ``vender 30 MWh para o submercado S em ago/2025 a 210 R\$/MWh''.
\end{itemize}

\subsection{Conjuntos do modelo}
\begin{itemize}
    \item $\mathcal{S}$: submercados.
    \item $\mathcal{T}^{F}$: meses futuros (os meses $t \ge t_0$ onde decidimos trades).
    \item $\mathcal{A}$: lista de trades disponíveis (cada trade $a$ tem um submercado $s(a)$ e mês $t(a)$).
    \item $\Omega$: cenários de preço ($|\Omega|$ grande;).
\end{itemize}

\subsection{Parâmetros (dados conhecidos)}
\begin{itemize}
    \item $G_{s,t}$: produção prevista (versão simples) no submercado $s$ e mês $t$.
          (Na versão completa, pode virar $G^\omega_{s,t}$.)
    \item $Q^{0,B}_{s,t}$ e $K^{0,B}_{s,t}$: compras já existentes no portfólio (volume e preço fixo).
    \item $Q^{0,S}_{s,t}$ e $K^{0,S}_{s,t}$: vendas já existentes no portfólio (volume e preço fixo).
    \item $K_a^{B}$: preço fixo (R\$/MWh) se o trade $a$ for \textbf{compra}.
    \item $K_a^{S}$: preço fixo (R\$/MWh) se o trade $a$ for \textbf{venda}.
    \item $\overline{q}_a^{B}$: limite de compra do trade $a$ (MWh).
    \item $\overline{q}_a^{S}$: limite de venda do trade $a$ (MWh).
    \item $P^\omega_{s,t}$: PLD no cenário $\omega$ para submercado $s$ e mês $t$ (construído na seção anterior).
    \item $\pi_\omega$: probabilidade do cenário.
    \item $\alpha \in (0,1)$: nível do CVaR (ex.: $0{,}95$).
    \item $\lambda \ge 0$: peso do risco (quanto maior, mais conservador).
\end{itemize}

\subsection{Variáveis de decisão (compra e venda separadas)}
Para cada trade $a\in\mathcal{A}$, $\omega \in \Omega$:
\[
q_{\omega a}^{B} \ge 0 \quad \text{(quanto comprar)},
\qquad
q_{\omega a}^{S} \ge 0 \quad \text{(quanto vender)}.
\]
Essas duas variáveis são didáticas: representam `compra'' e ``venda''.

\subsection{Como o trade entra na soma de compras e vendas por mês}
Defina o volume adicional comprado e vendido no mês $t$ e submercado $s$:
\[
Q^{\omega B}_{s,t} = \sum_{a\in\mathcal{A}:\, s(a)=s,\, t(a)=t} q_{\omega a}^{B},
\qquad
Q^{\omega S}_{s,t} = \sum_{a\in\mathcal{A}:\, s(a)=s,\, t(a)=t} q_{\omega a}^{S}.
\]

\subsection{Exposição (diferença) ao PLD por mês}
No mês $t$ e submercado $s$, a exposição (o ``resto'' que vai ao PLD) fica:
\[
E^{\omega}_{s,t} =
G^{\omega}_{s,t} + Q^{\omega 0 B}_{s t} + Q^{\omega B}_{s t}
-
\left(Q^{\omega 0 S}_{s t} + Q^{\omega S}_{s t}\right).
\]
Se $E^{\omega}_{s,t}<0$, faltou energia (compra no PLD). Se $E^{\omega}_{s,t}>0$, sobrou energia (vende no PLD).

\subsection{Lucro em um cenário}
O lucro total no cenário $\omega$ é:
\begin{align}
R^\omega
&=
\underbrace{
\sum_{s\in\mathcal{S}}\sum_{t\in\mathcal{T}^{F}}
\Big(K^{\omega 0 S}_{s t}Q^{\omega 0 S}_{s t} - K^{\omega 0 B}_{s t}Q^{\omega 0 B}_{s t}\Big)
}_{\text{contratos já existentes (constante)}}
\\[2mm]
&\quad+
\underbrace{
\sum_{a\in\mathcal{A}}\Big(K_{\omega a}^{S} q_{\omega a}^{S} - K_{\omega a}^{B} q_{\omega a}^{B}\Big)
}_{\text{novos contratos decididos (parte determinística)}}
\\[2mm]
&\quad+
\underbrace{
\sum_{s\in\mathcal{S}}\sum_{t\in\mathcal{T}^{F}}
E_{s,t}\cdot P^\omega_{s,t}
}_{\text{liquidação ao PLD (parte incerta)}}
.
\end{align}

\subsection{Risco via CVaR (totalmente linear)}
Trabalhamos diretamente com os lucros (sem inverter sinal).
Usamos variáveis auxiliares:
\[
\eta \in \mathbb{R},
\qquad
\xi_\omega \ge 0 \;\;\forall \omega\in\Omega.
\]
Com as restrições:
\[
\xi_\omega \ge \eta - R^\omega,
\qquad
\xi_\omega \ge 0,
\qquad \forall \omega\in\Omega.
\]
O CVaR é definido como:
\[
\text{CVaR} = \eta - \frac{1}{1-\alpha}\sum_{\omega\in\Omega}\pi_\omega \xi_\omega.
\]
Interpretação: $\eta$ é o quantil $(1-\alpha)$ dos lucros (VaR), e o CVaR representa o lucro médio nos $(1-\alpha)\%$ piores cenários. Quanto maior o CVaR, melhor (mais lucro nos cenários ruins).

\subsection{Função objetivo}
Maximizamos lucro esperado e maximizamos lucro nos piores cenários:
\[
\max \;\;
\sum_{\omega\in\Omega}\pi_\omega R^\omega
+
\lambda\left(
\eta - \frac{1}{1-\alpha}\sum_{\omega\in\Omega}\pi_\omega \xi_\omega
\right).
\]

\subsection{Restrições de volume (lineares)}
\[
0 \le q_{\omega a}^{B} \le \overline{q}_{\omega a}^{B},
\qquad
0 \le q_{\omega a}^{S} \le \overline{q}_{\omega a}^{S},
\qquad \forall a\in\mathcal{A}.
\]

\end{document}
